\begin{abstract}
Instruction-guided CAD editing requires changing a requested feature while preserving unrelated geometry. We study a frozen Astra model that produces original-coordinate program patches, combined with experience retrieval, graph search, execution feedback and patch repair. Candidates are checked for executable geometry, instruction-derived values, rendered-view semantic consistency, conditional finite-element response and robustness. On a locally defined 124-task BenchCAD holdout, the single Astra answer reaches \pipeAstra{} strict reference matches. The recorded full graph pipeline reaches \pipeFull{}, with \pipeFullGain{} recoveries and \pipeFullLoss{} regressions; it does not improve agreement. Retrieval alone reaches \pipeRag{}, but this configuration was identified after observing results and its paired gain is nonsignificant. Of \pipeWrong{} incorrect initial answers, \pipeWrongPass{} pass the geometry, numeric and FEM checks, exposing a gap between target-free feasibility and instruction correctness. The study reports a reproducible patch-repair workflow and its observed limitations. No model is trained in this pipeline, and physical checks use explicit diagnostic scenarios rather than inferred operating conditions.
\end{abstract}
